\section{Future work}

We plan to compare Representation Alignment across three single-domain datasets that span an explicit \emph{spatial}
difficulty axis: human faces (CelebA), multi-region architectural scenes (LSUN Church), and rigid objects under large
pose variation (Comprehensive Cars). CelebA is spatially homogeneous -- the face is centred and at a consistent scale,
so region separation is a low-effort task (Section~\ref{sec:celeba}) that may understate REPA's benefit -- whereas
Church and CompCars add multi-object scenes, perspective geometry, and viewpoint variation. This adjudicates between the
two competing hypotheses of Section~\ref{sec:task_setting}: that harder domains leave more for the teacher to contribute
(larger gains), or that the teacher's patch features are locally mismatched there (smaller or reversed gains). We sharpen
this into a falsifiable prediction: since DINOv2 is trained on web-scale data where buildings and vehicles are heavily
represented but tightly cropped faces are comparatively out of distribution, we expect its patch features to be more
spatially informative on Church and CompCars, and thus predict larger, more localized REPA gains there than on CelebA.

To keep ``domain'' from being entangled with ``dataset scale'', we will hold dataset size, compute, and teacher fixed
across domains, and report FID/KID speedups, the gradient-compatibility curves of Section~\ref{sec:celeba}, and
per-domain PCA feature maps. Finally, since REPA-$\Sigma$ (Section~\ref{sec:repa_sigma}) only acts when the alignment
and diffusion gradients conflict, we will test whether this low-noise conflict -- and hence the gain from gradient
surgery -- is itself domain-dependent, hypothesising it is strongest in the harder, more heterogeneous domains.
